% Figure: height-function stencil and GHF points (Chapter 4). Fractions,
% column heights, GHF points, weights and the fitted parabola:
% generated/heights.tex (generate.jl circle fixture); labels hand-written.
\begin{figure}
  \centering
  \begin{subfigure}[b]{0.38\textwidth}
    \centering
    \begin{tikzpicture}[x=0.56cm,y=0.56cm]
      \def\fxhpoint#1#2{}\def\fxgpoint#1#2#3#4#5#6{}\def\fxGcurve{}\def\fxGseg{}%
      % Pass 1: cells shaded by the fraction, the three columns outlined.
      \def\fxhcell#1#2#3#4#5{%
        \ifnum\pdfstrcmp{#1}{a}=0
          \pgfmathtruncatemacro\fxpct{100*#4}%
          \fill[trlight!\fxpct!white] (#2,#3) rectangle ++(1,1);
          \draw[tr cell] (#2,#3) rectangle ++(1,1);
          \ifnum#5>0 \node[tr tiny] at (#2+0.5,#3+0.5) {#4};\fi
        \fi}%
      % Generated by docs/techref/figures/generate.jl; do not edit.
\gdef\fxCircX{7.15}\gdef\fxCircY{6.35}\gdef\fxCircR{4.3}
\fxhcell{a}{6}{6}{1.00}{0}
\fxhcell{a}{7}{6}{1.00}{0}
\fxhcell{a}{8}{6}{1.00}{0}
\fxhcell{a}{9}{6}{1.00}{0}
\fxhcell{a}{10}{6}{1.00}{0}
\fxhcell{a}{6}{7}{1.00}{0}
\fxhcell{a}{7}{7}{1.00}{1}
\fxhcell{a}{8}{7}{1.00}{1}
\fxhcell{a}{9}{7}{1.00}{1}
\fxhcell{a}{10}{7}{1.00}{0}
\fxhcell{a}{6}{8}{1.00}{0}
\fxhcell{a}{7}{8}{1.00}{1}
\fxhcell{a}{8}{8}{1.00}{1}
\fxhcell{a}{9}{8}{1.00}{1}
\fxhcell{a}{10}{8}{0.84}{0}
\fxhcell{a}{6}{9}{1.00}{0}
\fxhcell{a}{7}{9}{1.00}{1}
\fxhcell{a}{8}{9}{1.00}{1}
\fxhcell{a}{9}{9}{0.88}{1}
\fxhcell{a}{10}{9}{0.16}{0}
\fxhcell{a}{6}{10}{0.59}{0}
\fxhcell{a}{7}{10}{0.63}{1}
\fxhcell{a}{8}{10}{0.42}{2}
\fxhcell{a}{9}{10}{0.05}{1}
\fxhcell{a}{10}{10}{0.00}{0}
\fxhcell{a}{6}{11}{0.00}{0}
\fxhcell{a}{7}{11}{0.00}{1}
\fxhcell{a}{8}{11}{0.00}{1}
\fxhcell{a}{9}{11}{0.00}{1}
\fxhcell{a}{10}{11}{0.00}{0}
\fxhcell{a}{6}{12}{0.00}{0}
\fxhcell{a}{7}{12}{0.00}{1}
\fxhcell{a}{8}{12}{0.00}{1}
\fxhcell{a}{9}{12}{0.00}{1}
\fxhcell{a}{10}{12}{0.00}{0}
\fxhcell{a}{6}{13}{0.00}{0}
\fxhcell{a}{7}{13}{0.00}{1}
\fxhcell{a}{8}{13}{0.00}{1}
\fxhcell{a}{9}{13}{0.00}{1}
\fxhcell{a}{10}{13}{0.00}{0}
\fxhpoint{7.5}{10.625924}
\fxhpoint{8.5}{10.421223}
\fxhpoint{9.5}{9.93439}
\gdef\fxHnx{0.326784}\gdef\fxHny{0.945099}\gdef\fxHi{9}\gdef\fxHj{11}
\fxhcell{b}{8}{7}{1.00}{0}
\fxhcell{b}{9}{7}{1.00}{0}
\fxhcell{b}{10}{7}{1.00}{0}
\fxhcell{b}{11}{7}{0.28}{0}
\fxhcell{b}{12}{7}{0.00}{0}
\fxhcell{b}{8}{8}{1.00}{0}
\fxhcell{b}{9}{8}{1.00}{1}
\fxhcell{b}{10}{8}{0.84}{1}
\fxhcell{b}{11}{8}{0.02}{1}
\fxhcell{b}{12}{8}{0.00}{0}
\fxhcell{b}{8}{9}{1.00}{0}
\fxhcell{b}{9}{9}{0.88}{1}
\fxhcell{b}{10}{9}{0.16}{2}
\fxhcell{b}{11}{9}{0.00}{1}
\fxhcell{b}{12}{9}{0.00}{0}
\fxhcell{b}{8}{10}{0.42}{0}
\fxhcell{b}{9}{10}{0.05}{1}
\fxhcell{b}{10}{10}{0.00}{1}
\fxhcell{b}{11}{10}{0.00}{1}
\fxhcell{b}{12}{10}{0.00}{0}
\fxhcell{b}{8}{11}{0.00}{0}
\fxhcell{b}{9}{11}{0.00}{0}
\fxhcell{b}{10}{11}{0.00}{0}
\fxhcell{b}{11}{11}{0.00}{0}
\fxhcell{b}{12}{11}{0.00}{0}
\fxgpoint{10.858881}{8.5}{1}{1}{0.60}{1}
\fxgpoint{9.5}{9.93439}{2}{1}{0.40}{1}
\fxgpoint{10.5}{9.004527}{2}{1}{0.40}{1}
\fxgpoint{10.045324}{9.5}{1}{1}{0.60}{1}
\fxgpoint{10.858881}{8.5}{1}{1}{0.60}{0}
\fxgpoint{10.858881}{8.5}{1}{1}{0.60}{0}
\fxgpoint{9.5}{9.93439}{2}{1}{0.40}{0}
\fxgpoint{10.045324}{9.5}{1}{1}{0.60}{0}
\fxgpoint{10.5}{9.004527}{2}{1}{0.40}{0}
\fxgpoint{10.045324}{9.5}{1}{1}{0.60}{0}
\fxgpoint{9.5}{9.93439}{2}{1}{0.40}{0}
\fxgpoint{10.5}{9.004527}{2}{1}{0.40}{0}
\gdef\fxGcurve{(11.137001,7.969132) -- (11.093626,8.062758) -- (11.048575,8.15502) -- (11.001847,8.245918) -- (10.953443,8.335452) -- (10.903363,8.423623) -- (10.851606,8.510429) -- (10.798172,8.595872) -- (10.743063,8.67995) -- (10.686276,8.762665) -- (10.627813,8.844016) -- (10.567674,8.924003) -- (10.505859,9.002626) -- (10.442367,9.079886) -- (10.377198,9.155781) -- (10.310353,9.230313) -- (10.241832,9.30348) -- (10.171634,9.375284) -- (10.099759,9.445724) -- (10.026209,9.5148) -- (9.950981,9.582512) -- (9.874078,9.64886) -- (9.795497,9.713844) -- (9.715241,9.777465) -- (9.633308,9.839721) -- (9.549698,9.900614) -- (9.464412,9.960143) -- (9.37745,10.018308) -- (9.288811,10.075109) -- (9.198496,10.130546) -- (9.106504,10.184619) -- (9.012836,10.237328) -- (8.917491,10.288674)}
\gdef\fxGox{10.256817}\gdef\fxGoy{9.315672}\gdef\fxGnx{0.775713}\gdef\fxGny{0.631086}\gdef\fxGi{11}\gdef\fxGj{10}
\gdef\fxGseg{(10.513633,9) -- (10,9.631343)}

      \draw[tr block] (\fxHi-2,\fxHj-4) rectangle (\fxHi+1,\fxHj+3);
      \draw[tr block,line width=1.2pt] (\fxHi-1,\fxHj-1) rectangle ++(1,1);
      % Pass 2: the column heights and the LHF normal.
      \def\fxhcell#1#2#3#4#5{}%
      \newcount\fxhk \fxhk=-2
      \def\fxhpoint#1#2{\global\advance\fxhk by 1
        \draw[trred,line width=0.6pt] (#1-0.5,#2) -- (#1+0.5,#2);
        \node[tr height] (h\the\fxhk) at (#1,#2) {};}%
      % Generated by docs/techref/figures/generate.jl; do not edit.
\gdef\fxCircX{7.15}\gdef\fxCircY{6.35}\gdef\fxCircR{4.3}
\fxhcell{a}{6}{6}{1.00}{0}
\fxhcell{a}{7}{6}{1.00}{0}
\fxhcell{a}{8}{6}{1.00}{0}
\fxhcell{a}{9}{6}{1.00}{0}
\fxhcell{a}{10}{6}{1.00}{0}
\fxhcell{a}{6}{7}{1.00}{0}
\fxhcell{a}{7}{7}{1.00}{1}
\fxhcell{a}{8}{7}{1.00}{1}
\fxhcell{a}{9}{7}{1.00}{1}
\fxhcell{a}{10}{7}{1.00}{0}
\fxhcell{a}{6}{8}{1.00}{0}
\fxhcell{a}{7}{8}{1.00}{1}
\fxhcell{a}{8}{8}{1.00}{1}
\fxhcell{a}{9}{8}{1.00}{1}
\fxhcell{a}{10}{8}{0.84}{0}
\fxhcell{a}{6}{9}{1.00}{0}
\fxhcell{a}{7}{9}{1.00}{1}
\fxhcell{a}{8}{9}{1.00}{1}
\fxhcell{a}{9}{9}{0.88}{1}
\fxhcell{a}{10}{9}{0.16}{0}
\fxhcell{a}{6}{10}{0.59}{0}
\fxhcell{a}{7}{10}{0.63}{1}
\fxhcell{a}{8}{10}{0.42}{2}
\fxhcell{a}{9}{10}{0.05}{1}
\fxhcell{a}{10}{10}{0.00}{0}
\fxhcell{a}{6}{11}{0.00}{0}
\fxhcell{a}{7}{11}{0.00}{1}
\fxhcell{a}{8}{11}{0.00}{1}
\fxhcell{a}{9}{11}{0.00}{1}
\fxhcell{a}{10}{11}{0.00}{0}
\fxhcell{a}{6}{12}{0.00}{0}
\fxhcell{a}{7}{12}{0.00}{1}
\fxhcell{a}{8}{12}{0.00}{1}
\fxhcell{a}{9}{12}{0.00}{1}
\fxhcell{a}{10}{12}{0.00}{0}
\fxhcell{a}{6}{13}{0.00}{0}
\fxhcell{a}{7}{13}{0.00}{1}
\fxhcell{a}{8}{13}{0.00}{1}
\fxhcell{a}{9}{13}{0.00}{1}
\fxhcell{a}{10}{13}{0.00}{0}
\fxhpoint{7.5}{10.625924}
\fxhpoint{8.5}{10.421223}
\fxhpoint{9.5}{9.93439}
\gdef\fxHnx{0.326784}\gdef\fxHny{0.945099}\gdef\fxHi{9}\gdef\fxHj{11}
\fxhcell{b}{8}{7}{1.00}{0}
\fxhcell{b}{9}{7}{1.00}{0}
\fxhcell{b}{10}{7}{1.00}{0}
\fxhcell{b}{11}{7}{0.28}{0}
\fxhcell{b}{12}{7}{0.00}{0}
\fxhcell{b}{8}{8}{1.00}{0}
\fxhcell{b}{9}{8}{1.00}{1}
\fxhcell{b}{10}{8}{0.84}{1}
\fxhcell{b}{11}{8}{0.02}{1}
\fxhcell{b}{12}{8}{0.00}{0}
\fxhcell{b}{8}{9}{1.00}{0}
\fxhcell{b}{9}{9}{0.88}{1}
\fxhcell{b}{10}{9}{0.16}{2}
\fxhcell{b}{11}{9}{0.00}{1}
\fxhcell{b}{12}{9}{0.00}{0}
\fxhcell{b}{8}{10}{0.42}{0}
\fxhcell{b}{9}{10}{0.05}{1}
\fxhcell{b}{10}{10}{0.00}{1}
\fxhcell{b}{11}{10}{0.00}{1}
\fxhcell{b}{12}{10}{0.00}{0}
\fxhcell{b}{8}{11}{0.00}{0}
\fxhcell{b}{9}{11}{0.00}{0}
\fxhcell{b}{10}{11}{0.00}{0}
\fxhcell{b}{11}{11}{0.00}{0}
\fxhcell{b}{12}{11}{0.00}{0}
\fxgpoint{10.858881}{8.5}{1}{1}{0.60}{1}
\fxgpoint{9.5}{9.93439}{2}{1}{0.40}{1}
\fxgpoint{10.5}{9.004527}{2}{1}{0.40}{1}
\fxgpoint{10.045324}{9.5}{1}{1}{0.60}{1}
\fxgpoint{10.858881}{8.5}{1}{1}{0.60}{0}
\fxgpoint{10.858881}{8.5}{1}{1}{0.60}{0}
\fxgpoint{9.5}{9.93439}{2}{1}{0.40}{0}
\fxgpoint{10.045324}{9.5}{1}{1}{0.60}{0}
\fxgpoint{10.5}{9.004527}{2}{1}{0.40}{0}
\fxgpoint{10.045324}{9.5}{1}{1}{0.60}{0}
\fxgpoint{9.5}{9.93439}{2}{1}{0.40}{0}
\fxgpoint{10.5}{9.004527}{2}{1}{0.40}{0}
\gdef\fxGcurve{(11.137001,7.969132) -- (11.093626,8.062758) -- (11.048575,8.15502) -- (11.001847,8.245918) -- (10.953443,8.335452) -- (10.903363,8.423623) -- (10.851606,8.510429) -- (10.798172,8.595872) -- (10.743063,8.67995) -- (10.686276,8.762665) -- (10.627813,8.844016) -- (10.567674,8.924003) -- (10.505859,9.002626) -- (10.442367,9.079886) -- (10.377198,9.155781) -- (10.310353,9.230313) -- (10.241832,9.30348) -- (10.171634,9.375284) -- (10.099759,9.445724) -- (10.026209,9.5148) -- (9.950981,9.582512) -- (9.874078,9.64886) -- (9.795497,9.713844) -- (9.715241,9.777465) -- (9.633308,9.839721) -- (9.549698,9.900614) -- (9.464412,9.960143) -- (9.37745,10.018308) -- (9.288811,10.075109) -- (9.198496,10.130546) -- (9.106504,10.184619) -- (9.012836,10.237328) -- (8.917491,10.288674)}
\gdef\fxGox{10.256817}\gdef\fxGoy{9.315672}\gdef\fxGnx{0.775713}\gdef\fxGny{0.631086}\gdef\fxGi{11}\gdef\fxGj{10}
\gdef\fxGseg{(10.513633,9) -- (10,9.631343)}

      \draw[tr vector] (h0) -- ++(1.6*\fxHnx,1.6*\fxHny) node[anchor=south,tr label] {$\inormal_c$};
      \node[tr small,anchor=north,fill=white,inner sep=1pt] at ([yshift=-2pt]h-1.south) {$\hdisc{-1}$};
      \node[tr small,anchor=north,fill=white,inner sep=1pt] at ([yshift=-2pt]h1.south) {$\hdisc{1}$};
      \node[tr small,anchor=north west,fill=white,inner sep=1pt] at ([xshift=1pt,yshift=-2pt]h0.south) {$\hdisc{0}$};
      \node[tr small,anchor=west] at (\fxHi+1.05,\fxHj-0.5) {$c$};
      \draw[tr thin vector] (\fxHi+1.6,\fxHj-3.6) -- (\fxHi+1.6,\fxHj-2.4)
        node[anchor=west,tr small] {$+\unitvec{2}$};
    \end{tikzpicture}
    \caption{LHF columns, principal direction $+\unitvec{2}$.}
    \label{fig:interface:heights:lhf}
  \end{subfigure}%
  \hfill
  \begin{subfigure}[b]{0.58\textwidth}
    \centering
    \begin{tikzpicture}[x=0.92cm,y=0.92cm]
      \def\fxhpoint#1#2{}\def\fxgpoint#1#2#3#4#5#6{}%
      \begin{scope}
        \clip (8,7) rectangle (13,12);
        \def\fxhcell#1#2#3#4#5{%
          \ifnum\pdfstrcmp{#1}{b}=0
            \pgfmathtruncatemacro\fxpct{100*#4}%
            \fill[trlight!\fxpct!white] (#2,#3) rectangle ++(1,1);
            \draw[tr cell] (#2,#3) rectangle ++(1,1);
            \node[tr tiny,anchor=south east,text=black!60] at (#2+1,#3) {#4};
          \fi}%
        % Generated by docs/techref/figures/generate.jl; do not edit.
\gdef\fxCircX{7.15}\gdef\fxCircY{6.35}\gdef\fxCircR{4.3}
\fxhcell{a}{6}{6}{1.00}{0}
\fxhcell{a}{7}{6}{1.00}{0}
\fxhcell{a}{8}{6}{1.00}{0}
\fxhcell{a}{9}{6}{1.00}{0}
\fxhcell{a}{10}{6}{1.00}{0}
\fxhcell{a}{6}{7}{1.00}{0}
\fxhcell{a}{7}{7}{1.00}{1}
\fxhcell{a}{8}{7}{1.00}{1}
\fxhcell{a}{9}{7}{1.00}{1}
\fxhcell{a}{10}{7}{1.00}{0}
\fxhcell{a}{6}{8}{1.00}{0}
\fxhcell{a}{7}{8}{1.00}{1}
\fxhcell{a}{8}{8}{1.00}{1}
\fxhcell{a}{9}{8}{1.00}{1}
\fxhcell{a}{10}{8}{0.84}{0}
\fxhcell{a}{6}{9}{1.00}{0}
\fxhcell{a}{7}{9}{1.00}{1}
\fxhcell{a}{8}{9}{1.00}{1}
\fxhcell{a}{9}{9}{0.88}{1}
\fxhcell{a}{10}{9}{0.16}{0}
\fxhcell{a}{6}{10}{0.59}{0}
\fxhcell{a}{7}{10}{0.63}{1}
\fxhcell{a}{8}{10}{0.42}{2}
\fxhcell{a}{9}{10}{0.05}{1}
\fxhcell{a}{10}{10}{0.00}{0}
\fxhcell{a}{6}{11}{0.00}{0}
\fxhcell{a}{7}{11}{0.00}{1}
\fxhcell{a}{8}{11}{0.00}{1}
\fxhcell{a}{9}{11}{0.00}{1}
\fxhcell{a}{10}{11}{0.00}{0}
\fxhcell{a}{6}{12}{0.00}{0}
\fxhcell{a}{7}{12}{0.00}{1}
\fxhcell{a}{8}{12}{0.00}{1}
\fxhcell{a}{9}{12}{0.00}{1}
\fxhcell{a}{10}{12}{0.00}{0}
\fxhcell{a}{6}{13}{0.00}{0}
\fxhcell{a}{7}{13}{0.00}{1}
\fxhcell{a}{8}{13}{0.00}{1}
\fxhcell{a}{9}{13}{0.00}{1}
\fxhcell{a}{10}{13}{0.00}{0}
\fxhpoint{7.5}{10.625924}
\fxhpoint{8.5}{10.421223}
\fxhpoint{9.5}{9.93439}
\gdef\fxHnx{0.326784}\gdef\fxHny{0.945099}\gdef\fxHi{9}\gdef\fxHj{11}
\fxhcell{b}{8}{7}{1.00}{0}
\fxhcell{b}{9}{7}{1.00}{0}
\fxhcell{b}{10}{7}{1.00}{0}
\fxhcell{b}{11}{7}{0.28}{0}
\fxhcell{b}{12}{7}{0.00}{0}
\fxhcell{b}{8}{8}{1.00}{0}
\fxhcell{b}{9}{8}{1.00}{1}
\fxhcell{b}{10}{8}{0.84}{1}
\fxhcell{b}{11}{8}{0.02}{1}
\fxhcell{b}{12}{8}{0.00}{0}
\fxhcell{b}{8}{9}{1.00}{0}
\fxhcell{b}{9}{9}{0.88}{1}
\fxhcell{b}{10}{9}{0.16}{2}
\fxhcell{b}{11}{9}{0.00}{1}
\fxhcell{b}{12}{9}{0.00}{0}
\fxhcell{b}{8}{10}{0.42}{0}
\fxhcell{b}{9}{10}{0.05}{1}
\fxhcell{b}{10}{10}{0.00}{1}
\fxhcell{b}{11}{10}{0.00}{1}
\fxhcell{b}{12}{10}{0.00}{0}
\fxhcell{b}{8}{11}{0.00}{0}
\fxhcell{b}{9}{11}{0.00}{0}
\fxhcell{b}{10}{11}{0.00}{0}
\fxhcell{b}{11}{11}{0.00}{0}
\fxhcell{b}{12}{11}{0.00}{0}
\fxgpoint{10.858881}{8.5}{1}{1}{0.60}{1}
\fxgpoint{9.5}{9.93439}{2}{1}{0.40}{1}
\fxgpoint{10.5}{9.004527}{2}{1}{0.40}{1}
\fxgpoint{10.045324}{9.5}{1}{1}{0.60}{1}
\fxgpoint{10.858881}{8.5}{1}{1}{0.60}{0}
\fxgpoint{10.858881}{8.5}{1}{1}{0.60}{0}
\fxgpoint{9.5}{9.93439}{2}{1}{0.40}{0}
\fxgpoint{10.045324}{9.5}{1}{1}{0.60}{0}
\fxgpoint{10.5}{9.004527}{2}{1}{0.40}{0}
\fxgpoint{10.045324}{9.5}{1}{1}{0.60}{0}
\fxgpoint{9.5}{9.93439}{2}{1}{0.40}{0}
\fxgpoint{10.5}{9.004527}{2}{1}{0.40}{0}
\gdef\fxGcurve{(11.137001,7.969132) -- (11.093626,8.062758) -- (11.048575,8.15502) -- (11.001847,8.245918) -- (10.953443,8.335452) -- (10.903363,8.423623) -- (10.851606,8.510429) -- (10.798172,8.595872) -- (10.743063,8.67995) -- (10.686276,8.762665) -- (10.627813,8.844016) -- (10.567674,8.924003) -- (10.505859,9.002626) -- (10.442367,9.079886) -- (10.377198,9.155781) -- (10.310353,9.230313) -- (10.241832,9.30348) -- (10.171634,9.375284) -- (10.099759,9.445724) -- (10.026209,9.5148) -- (9.950981,9.582512) -- (9.874078,9.64886) -- (9.795497,9.713844) -- (9.715241,9.777465) -- (9.633308,9.839721) -- (9.549698,9.900614) -- (9.464412,9.960143) -- (9.37745,10.018308) -- (9.288811,10.075109) -- (9.198496,10.130546) -- (9.106504,10.184619) -- (9.012836,10.237328) -- (8.917491,10.288674)}
\gdef\fxGox{10.256817}\gdef\fxGoy{9.315672}\gdef\fxGnx{0.775713}\gdef\fxGny{0.631086}\gdef\fxGi{11}\gdef\fxGj{10}
\gdef\fxGseg{(10.513633,9) -- (10,9.631343)}

        \draw[tr exact] (\fxCircX,\fxCircY) circle (\fxCircR);
        \draw[tr fit] \fxGcurve;
      \end{scope}
      \draw[tr block] (\fxGi-2,\fxGj-2) rectangle (\fxGi+1,\fxGj+1);
      \draw[tr block,line width=1.2pt] (\fxGi-1,\fxGj-1) rectangle ++(1,1);
      \draw[tr plic] \fxGseg;
      % GHF points: diamonds from x-columns, discs from y-columns; hollow:
      % discarded (zero weight or closer than h/2 to an earlier point).
      \def\fxhcell#1#2#3#4#5{}%
      \def\fxgpoint#1#2#3#4#5#6{\ifnum#6=0 \node[tr ghf rejected] at (#1,#2) {};\fi}%
      % Generated by docs/techref/figures/generate.jl; do not edit.
\gdef\fxCircX{7.15}\gdef\fxCircY{6.35}\gdef\fxCircR{4.3}
\fxhcell{a}{6}{6}{1.00}{0}
\fxhcell{a}{7}{6}{1.00}{0}
\fxhcell{a}{8}{6}{1.00}{0}
\fxhcell{a}{9}{6}{1.00}{0}
\fxhcell{a}{10}{6}{1.00}{0}
\fxhcell{a}{6}{7}{1.00}{0}
\fxhcell{a}{7}{7}{1.00}{1}
\fxhcell{a}{8}{7}{1.00}{1}
\fxhcell{a}{9}{7}{1.00}{1}
\fxhcell{a}{10}{7}{1.00}{0}
\fxhcell{a}{6}{8}{1.00}{0}
\fxhcell{a}{7}{8}{1.00}{1}
\fxhcell{a}{8}{8}{1.00}{1}
\fxhcell{a}{9}{8}{1.00}{1}
\fxhcell{a}{10}{8}{0.84}{0}
\fxhcell{a}{6}{9}{1.00}{0}
\fxhcell{a}{7}{9}{1.00}{1}
\fxhcell{a}{8}{9}{1.00}{1}
\fxhcell{a}{9}{9}{0.88}{1}
\fxhcell{a}{10}{9}{0.16}{0}
\fxhcell{a}{6}{10}{0.59}{0}
\fxhcell{a}{7}{10}{0.63}{1}
\fxhcell{a}{8}{10}{0.42}{2}
\fxhcell{a}{9}{10}{0.05}{1}
\fxhcell{a}{10}{10}{0.00}{0}
\fxhcell{a}{6}{11}{0.00}{0}
\fxhcell{a}{7}{11}{0.00}{1}
\fxhcell{a}{8}{11}{0.00}{1}
\fxhcell{a}{9}{11}{0.00}{1}
\fxhcell{a}{10}{11}{0.00}{0}
\fxhcell{a}{6}{12}{0.00}{0}
\fxhcell{a}{7}{12}{0.00}{1}
\fxhcell{a}{8}{12}{0.00}{1}
\fxhcell{a}{9}{12}{0.00}{1}
\fxhcell{a}{10}{12}{0.00}{0}
\fxhcell{a}{6}{13}{0.00}{0}
\fxhcell{a}{7}{13}{0.00}{1}
\fxhcell{a}{8}{13}{0.00}{1}
\fxhcell{a}{9}{13}{0.00}{1}
\fxhcell{a}{10}{13}{0.00}{0}
\fxhpoint{7.5}{10.625924}
\fxhpoint{8.5}{10.421223}
\fxhpoint{9.5}{9.93439}
\gdef\fxHnx{0.326784}\gdef\fxHny{0.945099}\gdef\fxHi{9}\gdef\fxHj{11}
\fxhcell{b}{8}{7}{1.00}{0}
\fxhcell{b}{9}{7}{1.00}{0}
\fxhcell{b}{10}{7}{1.00}{0}
\fxhcell{b}{11}{7}{0.28}{0}
\fxhcell{b}{12}{7}{0.00}{0}
\fxhcell{b}{8}{8}{1.00}{0}
\fxhcell{b}{9}{8}{1.00}{1}
\fxhcell{b}{10}{8}{0.84}{1}
\fxhcell{b}{11}{8}{0.02}{1}
\fxhcell{b}{12}{8}{0.00}{0}
\fxhcell{b}{8}{9}{1.00}{0}
\fxhcell{b}{9}{9}{0.88}{1}
\fxhcell{b}{10}{9}{0.16}{2}
\fxhcell{b}{11}{9}{0.00}{1}
\fxhcell{b}{12}{9}{0.00}{0}
\fxhcell{b}{8}{10}{0.42}{0}
\fxhcell{b}{9}{10}{0.05}{1}
\fxhcell{b}{10}{10}{0.00}{1}
\fxhcell{b}{11}{10}{0.00}{1}
\fxhcell{b}{12}{10}{0.00}{0}
\fxhcell{b}{8}{11}{0.00}{0}
\fxhcell{b}{9}{11}{0.00}{0}
\fxhcell{b}{10}{11}{0.00}{0}
\fxhcell{b}{11}{11}{0.00}{0}
\fxhcell{b}{12}{11}{0.00}{0}
\fxgpoint{10.858881}{8.5}{1}{1}{0.60}{1}
\fxgpoint{9.5}{9.93439}{2}{1}{0.40}{1}
\fxgpoint{10.5}{9.004527}{2}{1}{0.40}{1}
\fxgpoint{10.045324}{9.5}{1}{1}{0.60}{1}
\fxgpoint{10.858881}{8.5}{1}{1}{0.60}{0}
\fxgpoint{10.858881}{8.5}{1}{1}{0.60}{0}
\fxgpoint{9.5}{9.93439}{2}{1}{0.40}{0}
\fxgpoint{10.045324}{9.5}{1}{1}{0.60}{0}
\fxgpoint{10.5}{9.004527}{2}{1}{0.40}{0}
\fxgpoint{10.045324}{9.5}{1}{1}{0.60}{0}
\fxgpoint{9.5}{9.93439}{2}{1}{0.40}{0}
\fxgpoint{10.5}{9.004527}{2}{1}{0.40}{0}
\gdef\fxGcurve{(11.137001,7.969132) -- (11.093626,8.062758) -- (11.048575,8.15502) -- (11.001847,8.245918) -- (10.953443,8.335452) -- (10.903363,8.423623) -- (10.851606,8.510429) -- (10.798172,8.595872) -- (10.743063,8.67995) -- (10.686276,8.762665) -- (10.627813,8.844016) -- (10.567674,8.924003) -- (10.505859,9.002626) -- (10.442367,9.079886) -- (10.377198,9.155781) -- (10.310353,9.230313) -- (10.241832,9.30348) -- (10.171634,9.375284) -- (10.099759,9.445724) -- (10.026209,9.5148) -- (9.950981,9.582512) -- (9.874078,9.64886) -- (9.795497,9.713844) -- (9.715241,9.777465) -- (9.633308,9.839721) -- (9.549698,9.900614) -- (9.464412,9.960143) -- (9.37745,10.018308) -- (9.288811,10.075109) -- (9.198496,10.130546) -- (9.106504,10.184619) -- (9.012836,10.237328) -- (8.917491,10.288674)}
\gdef\fxGox{10.256817}\gdef\fxGoy{9.315672}\gdef\fxGnx{0.775713}\gdef\fxGny{0.631086}\gdef\fxGi{11}\gdef\fxGj{10}
\gdef\fxGseg{(10.513633,9) -- (10,9.631343)}

      \def\fxgpoint#1#2#3#4#5#6{%
        \ifnum#6=1
          \ifnum#3=1 \node[tr ghf x] at (#1,#2) {};\else\node[tr ghf y] at (#1,#2) {};\fi
          \node[tr tiny,anchor=north east,inner sep=1.5pt,fill=white,fill opacity=0.8,text opacity=1] at (#1,#2) {$#5$};
        \fi}%
      % Generated by docs/techref/figures/generate.jl; do not edit.
\gdef\fxCircX{7.15}\gdef\fxCircY{6.35}\gdef\fxCircR{4.3}
\fxhcell{a}{6}{6}{1.00}{0}
\fxhcell{a}{7}{6}{1.00}{0}
\fxhcell{a}{8}{6}{1.00}{0}
\fxhcell{a}{9}{6}{1.00}{0}
\fxhcell{a}{10}{6}{1.00}{0}
\fxhcell{a}{6}{7}{1.00}{0}
\fxhcell{a}{7}{7}{1.00}{1}
\fxhcell{a}{8}{7}{1.00}{1}
\fxhcell{a}{9}{7}{1.00}{1}
\fxhcell{a}{10}{7}{1.00}{0}
\fxhcell{a}{6}{8}{1.00}{0}
\fxhcell{a}{7}{8}{1.00}{1}
\fxhcell{a}{8}{8}{1.00}{1}
\fxhcell{a}{9}{8}{1.00}{1}
\fxhcell{a}{10}{8}{0.84}{0}
\fxhcell{a}{6}{9}{1.00}{0}
\fxhcell{a}{7}{9}{1.00}{1}
\fxhcell{a}{8}{9}{1.00}{1}
\fxhcell{a}{9}{9}{0.88}{1}
\fxhcell{a}{10}{9}{0.16}{0}
\fxhcell{a}{6}{10}{0.59}{0}
\fxhcell{a}{7}{10}{0.63}{1}
\fxhcell{a}{8}{10}{0.42}{2}
\fxhcell{a}{9}{10}{0.05}{1}
\fxhcell{a}{10}{10}{0.00}{0}
\fxhcell{a}{6}{11}{0.00}{0}
\fxhcell{a}{7}{11}{0.00}{1}
\fxhcell{a}{8}{11}{0.00}{1}
\fxhcell{a}{9}{11}{0.00}{1}
\fxhcell{a}{10}{11}{0.00}{0}
\fxhcell{a}{6}{12}{0.00}{0}
\fxhcell{a}{7}{12}{0.00}{1}
\fxhcell{a}{8}{12}{0.00}{1}
\fxhcell{a}{9}{12}{0.00}{1}
\fxhcell{a}{10}{12}{0.00}{0}
\fxhcell{a}{6}{13}{0.00}{0}
\fxhcell{a}{7}{13}{0.00}{1}
\fxhcell{a}{8}{13}{0.00}{1}
\fxhcell{a}{9}{13}{0.00}{1}
\fxhcell{a}{10}{13}{0.00}{0}
\fxhpoint{7.5}{10.625924}
\fxhpoint{8.5}{10.421223}
\fxhpoint{9.5}{9.93439}
\gdef\fxHnx{0.326784}\gdef\fxHny{0.945099}\gdef\fxHi{9}\gdef\fxHj{11}
\fxhcell{b}{8}{7}{1.00}{0}
\fxhcell{b}{9}{7}{1.00}{0}
\fxhcell{b}{10}{7}{1.00}{0}
\fxhcell{b}{11}{7}{0.28}{0}
\fxhcell{b}{12}{7}{0.00}{0}
\fxhcell{b}{8}{8}{1.00}{0}
\fxhcell{b}{9}{8}{1.00}{1}
\fxhcell{b}{10}{8}{0.84}{1}
\fxhcell{b}{11}{8}{0.02}{1}
\fxhcell{b}{12}{8}{0.00}{0}
\fxhcell{b}{8}{9}{1.00}{0}
\fxhcell{b}{9}{9}{0.88}{1}
\fxhcell{b}{10}{9}{0.16}{2}
\fxhcell{b}{11}{9}{0.00}{1}
\fxhcell{b}{12}{9}{0.00}{0}
\fxhcell{b}{8}{10}{0.42}{0}
\fxhcell{b}{9}{10}{0.05}{1}
\fxhcell{b}{10}{10}{0.00}{1}
\fxhcell{b}{11}{10}{0.00}{1}
\fxhcell{b}{12}{10}{0.00}{0}
\fxhcell{b}{8}{11}{0.00}{0}
\fxhcell{b}{9}{11}{0.00}{0}
\fxhcell{b}{10}{11}{0.00}{0}
\fxhcell{b}{11}{11}{0.00}{0}
\fxhcell{b}{12}{11}{0.00}{0}
\fxgpoint{10.858881}{8.5}{1}{1}{0.60}{1}
\fxgpoint{9.5}{9.93439}{2}{1}{0.40}{1}
\fxgpoint{10.5}{9.004527}{2}{1}{0.40}{1}
\fxgpoint{10.045324}{9.5}{1}{1}{0.60}{1}
\fxgpoint{10.858881}{8.5}{1}{1}{0.60}{0}
\fxgpoint{10.858881}{8.5}{1}{1}{0.60}{0}
\fxgpoint{9.5}{9.93439}{2}{1}{0.40}{0}
\fxgpoint{10.045324}{9.5}{1}{1}{0.60}{0}
\fxgpoint{10.5}{9.004527}{2}{1}{0.40}{0}
\fxgpoint{10.045324}{9.5}{1}{1}{0.60}{0}
\fxgpoint{9.5}{9.93439}{2}{1}{0.40}{0}
\fxgpoint{10.5}{9.004527}{2}{1}{0.40}{0}
\gdef\fxGcurve{(11.137001,7.969132) -- (11.093626,8.062758) -- (11.048575,8.15502) -- (11.001847,8.245918) -- (10.953443,8.335452) -- (10.903363,8.423623) -- (10.851606,8.510429) -- (10.798172,8.595872) -- (10.743063,8.67995) -- (10.686276,8.762665) -- (10.627813,8.844016) -- (10.567674,8.924003) -- (10.505859,9.002626) -- (10.442367,9.079886) -- (10.377198,9.155781) -- (10.310353,9.230313) -- (10.241832,9.30348) -- (10.171634,9.375284) -- (10.099759,9.445724) -- (10.026209,9.5148) -- (9.950981,9.582512) -- (9.874078,9.64886) -- (9.795497,9.713844) -- (9.715241,9.777465) -- (9.633308,9.839721) -- (9.549698,9.900614) -- (9.464412,9.960143) -- (9.37745,10.018308) -- (9.288811,10.075109) -- (9.198496,10.130546) -- (9.106504,10.184619) -- (9.012836,10.237328) -- (8.917491,10.288674)}
\gdef\fxGox{10.256817}\gdef\fxGoy{9.315672}\gdef\fxGnx{0.775713}\gdef\fxGny{0.631086}\gdef\fxGi{11}\gdef\fxGj{10}
\gdef\fxGseg{(10.513633,9) -- (10,9.631343)}

      \node[circle,fill=trpurple,inner sep=0pt,minimum size=3.6pt] (o) at (\fxGox,\fxGoy) {};
      \draw[tr vector] (o) -- ++(0.75*\fxGnx,0.75*\fxGny) node[anchor=south west,tr label] {$\inormal_c$};
      \draw[tr thin vector,black!70] (o) -- ++(-0.75*\fxGny,0.75*\fxGnx) node[anchor=south east,tr small] {$\itangent$};
      \node[tr small,anchor=north east,text=trpurple] at (o.south west) {$\icentroid{c}$};
      \node[tr small,anchor=north west] at (\fxGi-1,\fxGj-1) {$c$};
    \end{tikzpicture}
    \caption{GHF points of the $3\times3$ neighbourhood, weights $\fitweight{j}$, fit (purple).}
    \label{fig:interface:heights:ghf}
  \end{subfigure}
  \caption{Height functions on the circle fixture of \cref{fig:interface:plic}
    (shading: $\condfrac{\liq}{c}$). (a) The seven-cell columns through $c$ and
    its two neighbours give the heights $\hdisc{-1},\hdisc{0},\hdisc{1}$
    of~\cref{eq:interface:height} (red) and the LHF normal
    \cref{eq:interface:lhf}; no other signed direction is valid here. (b) A cell
    without valid LHF columns (W in \cref{fig:interface:plic:patch}): the heights
    of all signed directions in the outlined $3\times3$ block are candidate points
    $\ghfpoint{j}$; \chkGhfPoints\ are kept with the weights
    $\fitweight{j}=\max(0,\pm\inormal_c\cdot\unitvec{k})^2$ of
    \cref{eq:interface:ghf}, \chkGhfRejected\ are discarded (hollow). The parabola
    in the frame $(\inormal_c,\itangent)$ at $\icentroid{c}$ of the (WY) plane gives
    the curvature. Dashed: the exact circle.}
  \label{fig:interface:heights}
\end{figure}
